% chapters/05_rl_convergence.tex

\section{Overview}

This chapter establishes convergence guarantees for a class of multi-timescale
actor-critic algorithms designed to learn the mean-field game equilibrium with common
noise without prior knowledge of the model dynamics. While Chapter 3 proved that the
equilibrium exists and is unique under appropriate regularity conditions, and Chapter
4 provided numerical methods to compute it when the model is fully known, the present
chapter addresses the setting where the drift $b$, the reward functions $f$ and $g$,
or the equilibrium measure $\mu^*$ are unknown and must be learned from observed
trajectories.

The central contribution is twofold. First, we prove that a carefully designed
three-timescale actor-critic algorithm converges along subsequences (in the sense of
$\liminf W_2 = 0$) to the MFG equilibrium measure flow in the Wasserstein-2 topology,
under strong regularity and compactness assumptions on the function approximators
(Theorem~\ref{thm:rl-subsequential}). The proof adapts \citet{borkar1997,borkar2008}'s (1997, 2008) ordinary
differential equation (ODE) method for stochastic approximation to the measure-valued
setting, leveraging the contraction property established in
Theorem~\ref{thm:wellposedness}. \textbf{We emphasize that the convergence is only
along subsequences; full almost-sure convergence would require additional uniqueness
of limit points, which is not proved here.}

Second, we show that when the policy is additionally regularised with an entropy
term, the convergence can be strengthened to an exponential rate, provided the policy
class satisfies a log-Sobolev inequality (Theorem~\ref{thm:rl-exponential}). While
this assumption is restrictive and serves as an idealised benchmark, it illustrates
the structural conditions under which fast convergence is achievable and provides a
theoretical bridge between entropy-regularised RL and mean-field games.

\paragraph{Prior work.} The application of reinforcement learning to MFGs has seen
rapid progress. \citep{angiulietal2022} introduced a unified
two-timescale RL algorithm for MFGs and mean-field control. \citep{hulauriere2023}
developed a two-timescale actor-critic method specifically for common-noise MFGs.
Most relevantly, \citep{angiulietal2024} proved convergence of a
three-timescale Q-learning algorithm for mean-field control games with common noise,
establishing that the multi-scale dynamics correctly capture the representative
agent's viewpoint. \textbf{The present chapter extends this ODE-method approach to
the continuous-time, continuous-space setting with explicit Wasserstein distance
control. The proof structure closely follows that of \citep{angiulietal2024}; the novel
technical challenge is establishing tightness and convergence of the measure-valued
component in the Wasserstein topology.} Other relevant works include the fictitious
play algorithms of Perrin et al.\ (2022) and the mean-field actor-critic flow of \citep{zhouzhouhu2025}. \Cref{tab:rl-comparison} situates our contributions relative to
this literature.

\begin{table}[ht]
\centering
\small
\caption{Comparison of RL for MFG Convergence Results}
\label{tab:rl-comparison}
\begin{tabular}{p{2.3cm}p{1.6cm}p{1.4cm}p{1.6cm}p{1.6cm}p{1.8cm}p{1.8cm}}
\toprule
\textbf{Work} & \textbf{Setting} & \textbf{Time-scales} & \textbf{Common Noise} &
    \textbf{Wasser-stein Control} & \textbf{Function Class} & \textbf{Convergence
    Type} \\
\midrule
Angiuli et al.\ (2022) & MFG/MFC & 2 & No & No & Finite & Almost-sure \\
Hu \& Lauri\`ere (2023) & MFG & 2 & Yes & No & Finite & Almost-sure \\
Angiuli et al.\ (2024) & MFC & 3 & Yes & No & Finite & Almost-sure \\
Zhou et al.\ (2025) & MFG & Continuous & No & Yes & Parametric & Weak convergence \\
\textbf{This thesis} (Thm~\ref{thm:rl-subsequential}) & MFG & 3 & Yes & Yes
    ($\liminf$) & Compact Lipschitz & Subsequential a.s. \\
\textbf{This thesis} (Thm~\ref{thm:rl-exponential}) & MFG & 3 & Yes & Yes
    (exponential) & LSI policy class & Exponential rate \\
\bottomrule
\end{tabular}
\end{table}

\paragraph{Scope of the guarantees.} \textbf{The convergence results in Theorems
\ref{thm:rl-subsequential} and \ref{thm:rl-exponential} are established for
regularised, compact function classes and under strong structural assumptions
(Lipschitz continuity, boundedness, stability of limiting ODEs, and log-Sobolev
inequality for the exponential rate). These assumptions are restrictive and are not
satisfied by arbitrary deep neural networks as commonly used in practice. The results
should be interpreted as theoretical baselines --- existence and consistency proofs
for a class of idealised algorithms --- rather than performance guarantees for
unconstrained deep reinforcement learning implementations. Extending these guarantees
to overparameterised networks remains an important open problem, discussed further in
Chapter 10.}

\paragraph{Chapter organisation.} \Cref{sec:rl-formulation} formulates the RL problem.
\Cref{sec:rl-architecture} presents the three-timescale architecture.
\Cref{sec:rl-theorems} states the main theorems. \Cref{sec:rl-proof-subsequential}
proves subsequential almost-sure convergence via the ODE method.
\Cref{sec:rl-entropy} analyses entropy regularisation and exponential convergence.
\Cref{sec:rl-implementation} discusses practical implementation and limitations.
\Cref{sec:rl-conclusion} concludes.


\section{Reinforcement Learning Formulation of the MFG}
\label{sec:rl-formulation}

We recall the mean-field game setting from Chapter 1. A representative agent controls
a state process $X_t \in \mathbb{R}^d$ evolving according to the McKean--Vlasov SDE
$dX_t = b(t,X_t,\mu_t,\alpha_t)\, dt + \sigma\, dW_t^i + \sigma_0\, dW_t^0$, $X_0 \sim
\mu_{\mathrm{init}}$, where $\alpha_t \in A \subset \mathbb{R}^k$ is the control
(action), and $\mu_t = \mathcal{L}(X_t \mid \mathcal{F}_t^{W^0})$ is the conditional
law of the state given the common noise. The agent aims to maximise the expected
cumulative reward
\[
    J(\alpha; \mu) = \mathbb{E}\left[ \int_0^T f(t,X_t,\mu_t,\alpha_t)\, dt +
    g(X_T,\mu_T) \right],
\]
taking the measure flow $\mu = (\mu_t)_{t \in [0,T]}$ as given. The MFG equilibrium is
a pair $(\mu^*, \alpha^*)$ such that $\alpha^*$ is optimal for the agent given $\mu^*$,
and $\mu_t^* = \mathcal{L}(X_t^* \mid \mathcal{F}_t^{W^0})$ where $X^*$ is the state
process under $\alpha^*$.

In the reinforcement learning setting, the agent does not have access to the
functional forms of $b$, $f$, or $g$. Instead, it observes a stream of data: at
discrete time steps $t_n = n\Delta t$, the agent observes the current state $X_n$,
takes an action $\alpha_n$, receives a reward
$R_n \approx f(t_n, X_n, \hat\mu_n, \alpha_n)\Delta t$, and observes the next state
$X_{n+1}$. The measure flow $\hat\mu_n$ must also be estimated, either from a large
population of agents or by maintaining an empirical distribution of historical
states. The goal is to learn a policy
$\pi : [0,T] \times \mathbb{R}^d \times \mathcal{P}_2(\mathbb{R}^d) \to A$ that
approximates the optimal feedback control $\alpha_t^* = \alpha^*(t, X_t, \mu_t)$.

By the Stochastic Maximum Principle (Theorem~\ref{thm:smp}), the optimal control is
characterised by the Hamiltonian maximisation condition
$\alpha_t^* = \argmax_{\alpha \in A} H(t, X_t, \mu_t, p_t, \alpha)$, where
$p_t = D_x u(t, X_t, \mu_t)$ is the adjoint process. Thus, if the agent can learn the
adjoint process $p_t$ (the critic) and the optimal policy (the actor), it can
approximate the equilibrium. The presence of common noise introduces a third
learning timescale: the measure flow $\mu_t$ evolves stochastically and must be
tracked by a separate estimator (the common-noise filter).


\section{Three-Timescale Actor-Critic Architecture}
\label{sec:rl-architecture}

We propose a three-timescale stochastic approximation algorithm consisting of the
following components:

\begin{enumerate}
    \item \textbf{Critic (fastest timescale)}: Estimates the adjoint process $p_t$.
    In the continuous-time limit, the critic solves the BSDE \eqref{eq:adjoint-bsde}.
    \item \textbf{Actor (intermediate timescale)}: Updates the policy parameters to
    maximise the Hamiltonian, using the critic's estimates.
    \item \textbf{Common-noise filter (slowest timescale)}: Tracks the population
    measure flow $\mu_t$, which evolves according to the Fokker--Planck SPDE
    \eqref{eq:fp-spde-weak}.
\end{enumerate}

The timescale separation is essential: the critic must adapt quickly to changes in
the policy and the measure; the actor updates more slowly, treating the critic's
value estimates as approximately converged; the measure filter updates most slowly,
reflecting the fact that the population distribution changes gradually relative to
individual decisions.

\subsection{Discrete-Time Algorithm}

Let $\gamma_t, \beta_t, \eta_t$ be three positive step-size sequences satisfying the
Robbins--Monro conditions
\[
    \sum_{t=0}^\infty \gamma_t = \sum_{t=0}^\infty \beta_t = \sum_{t=0}^\infty \eta_t =
    \infty, \qquad \sum_{t=0}^\infty (\gamma_t^2 + \beta_t^2 + \eta_t^2) < \infty,
\]
and the timescale separation condition
\begin{equation}
    \lim_{t \to \infty} \frac{\beta_t}{\gamma_t} = 0, \qquad \lim_{t \to \infty}
    \frac{\eta_t}{\beta_t} = 0.
    \label{eq:timescale-separation}
\end{equation}
Condition \eqref{eq:timescale-separation} ensures that, relative to the critic, the
actor and the filter evolve quasi-statically.

We parametrise the critic as $p_\theta(t,x,\mu)$, the actor as a policy
$\pi_\phi(t,x,\mu)$, and maintain an empirical measure $\hat\mu_t$ of the population
state. The algorithm proceeds as follows at each discrete time step $n$:

\begin{enumerate}
    \item Observe current state $X_n$, action $\alpha_n$, reward $R_n$, and next
    state $X_{n+1}$.
    \item \textbf{Critic update (fast)}:
    \[
        \theta_{n+1} = \theta_n + \gamma_n \delta_n \nabla_\theta p_{\theta_n}(t_n,
        X_n, \hat\mu_n),
    \]
    where $\delta_n$ is the temporal difference error for the adjoint process. In the
    continuous-time limit, $\delta_n$ approximates the BSDE residual:
    $\delta_n = p_{\theta_n}(t_{n+1}, X_{n+1}, \hat\mu_{n+1}) - p_{\theta_n}(t_n, X_n,
    \hat\mu_n) + \partial_x H(t_n, X_n, \hat\mu_n, p_{\theta_n}, \alpha_n) \Delta t$.
    \item \textbf{Actor update (intermediate)}:
    \[
        \phi_{n+1} = \phi_n + \beta_n \nabla_\phi H(t_n, X_n, \hat\mu_n, p_{\theta_n},
        \pi_{\phi_n}) \cdot \nabla_\phi \pi_{\phi_n}(t_n, X_n, \hat\mu_n).
    \]
    \item \textbf{Measure filter update (slow)}:
    \[
        \hat\mu_{n+1} = (1-\eta_n)\hat\mu_n + \eta_n \delta_{X_{n+1}}.
    \]
    In practice, a finite window of recent states is maintained, or a parametric
    density estimator is updated.
    \item \textbf{Action selection}: $\alpha_{n+1} = \pi_{\phi_{n+1}}(t_{n+1},
    X_{n+1}, \hat\mu_{n+1})$ (possibly with exploration noise).
\end{enumerate}

\subsection{Continuous-Time Limiting Dynamics}

Under the timescale separation \eqref{eq:timescale-separation}, the coupled
stochastic approximation can be analysed by considering three limiting ordinary
differential equations (ODEs). As $\gamma_n \to 0$ faster than $\beta_n, \eta_n$, the
critic parameters $\theta$ converge to the solution of a deterministic equation
parameterised by the slower variables $\phi$ and $\hat\mu$. In turn, the actor
parameters $\phi$ follow an ODE driven by the converged critic, and the measure
$\hat\mu$ follows the slowest dynamics, the Fokker--Planck SPDE.

Formally, let $\tau_\gamma, \tau_\beta, \tau_\eta$ denote the three timescales. Define
the limiting critic ODE:
\begin{equation}
    \dot\theta = \mathbb{E}_{X \sim \hat\mu}\big[ \delta(\theta, \phi, \hat\mu)
    \nabla_\theta p_\theta(X, \hat\mu) \big],
    \label{eq:critic-ode}
\end{equation}
the limiting actor ODE:
\begin{equation}
    \dot\phi = \mathbb{E}_{X \sim \hat\mu}\big[ \nabla_\phi H(X, \hat\mu,
    p_{\theta^*(\phi,\hat\mu)}, \pi_\phi) \cdot \nabla_\phi \pi_\phi(X, \hat\mu) \big],
    \label{eq:actor-ode}
\end{equation}
where $\theta^*(\phi, \hat\mu)$ is the equilibrium of \eqref{eq:critic-ode}, and the
measure evolution:
\begin{equation}
    d\hat\mu_t = \mathcal{L}_\phi^* \hat\mu_t\, dt + \sigma_0 \nabla \hat\mu_t\,
    dW_t^0,
    \label{eq:measure-filter-spde}
\end{equation}
where $\mathcal{L}_\phi^*$ is the Fokker--Planck operator under the policy $\pi_\phi$.
The full derivation of these limiting dynamics is provided in Appendix~A.


\section{Statement of Main Theorems}
\label{sec:rl-theorems}

We now state the principal convergence results for the three-timescale algorithm. The
proofs are developed in \Cref{sec:rl-proof-subsequential} and \Cref{sec:rl-entropy},
with detailed algebraic manipulations deferred to Appendix~A.

\begin{theorem}[Subsequential Almost-Sure Convergence of Multi-Timescale Actor-Critic]
\label{thm:rl-subsequential}
Let Assumptions 1.13--1.19 hold. Assume that the function approximators $p_\theta$
and $\pi_\phi$ belong to a compact parameter class $\Theta \times \Phi$ and are
uniformly bounded and Lipschitz continuous in their parameters and inputs. Suppose
the step-size sequences satisfy the Robbins--Monro conditions and the timescale
separation \eqref{eq:timescale-separation}. Assume further that the limiting ODEs
\eqref{eq:critic-ode}--\eqref{eq:actor-ode} possess globally asymptotically stable
equilibria $\theta^*(\phi,\hat\mu)$ and $\phi^*(\hat\mu)$, respectively. Then, for
almost all initialisations, the three-timescale algorithm converges
$\mathbb{P}$-almost surely along subsequences in the following sense:
\[
    \liminf_{t \to \infty} W_2(\hat\mu_t, \mu_t^*) = 0 \qquad \mathbb{P}\text{-a.s.},
\]
where $\mu^*$ is the unique MFG equilibrium measure flow. Moreover, along any
convergent subsequence, the policy parameters $\phi_n$ converge to a set of
parameters inducing the optimal feedback control $\alpha^*$.
\end{theorem}

\begin{remark}[Function approximator assumptions]
The assumptions of compactness and Lipschitz continuity are satisfied, for example, by
linear combinations of bounded basis functions (e.g., Fourier features, radial basis
functions) with coefficients constrained to a compact set, or by neural networks with
bounded weights and spectral normalisation. \textbf{They are not satisfied by
standard unconstrained deep neural networks.} The theorem provides a theoretical
foundation for idealised algorithms; extending the proof to overparameterised
networks is an open problem (see Chapter 10).
\end{remark}

\begin{remark}[Subsequential convergence]
The theorem establishes $\liminf W_2 = 0$, not full convergence. This is a weaker
result than almost-sure convergence to the unique equilibrium. The reason is that the
stochastic approximation argument only guarantees that the iterates visit
neighbourhoods of the equilibrium infinitely often; proving that the sequence
eventually stays near the equilibrium requires additional monotonicity or
uniqueness-of-limit-point arguments that are not available in the measure-valued
setting. This limitation should be kept in mind when interpreting the result.
\end{remark}

\begin{remark}[Stability of limiting ODEs]
The assumption that the limiting ODEs possess globally asymptotically stable
equilibria is a strong but standard condition in stochastic approximation theory
(\citep{borkar2008}, Chapter 2). In the LQ-MFG case, this stability can be verified
explicitly (see Appendix~A). For general nonlinear MFGs, it remains an active area
of research; recent work by \citep{zhouzhouhu2025} establishes such stability for
mean-field actor-critic flows under certain convexity conditions.
\end{remark}

\begin{theorem}[Exponential Convergence under Entropy Regularisation]
\label{thm:rl-exponential}
In addition to the hypotheses of Theorem~\ref{thm:rl-subsequential}, suppose that the
policy is entropy-regularised with coefficient $\beta > 0$, so that the objective
includes a term $-\beta\, \mathrm{KL}(\pi \| \pi_0)$. Assume further that the policy
class satisfies a log-Sobolev inequality with constant $C_{\mathrm{Lip}}$ (e.g.,
Gaussian policies or log-concave policy classes). Then the Wasserstein distance
between the learned measure flow $\hat\mu_t$ and the equilibrium $\mu^*$ satisfies
\[
    W_2(\hat\mu_t, \mu^*) \leq C_0 e^{-\lambda t} W_2(\hat\mu_0, \mu^*), \qquad
    \lambda = \frac{\beta}{2 C_{\mathrm{Lip}}},
\]
where $C_0$ depends on the initialisation and the model parameters. This
exponential rate holds for the continuous-time limiting dynamics and provides an
upper bound on the convergence speed of the discrete algorithm before discretisation
errors dominate.
\end{theorem}

\begin{remark}[Log-Sobolev assumption]
The log-Sobolev inequality is a strong structural condition that typically holds for
Gaussian policies or policies that are strongly log-concave in the parameters. In the
context of neural network policies, the LSI generally fails unless explicit
regularisation (such as weight decay or spectral normalisation) is imposed. The
theorem is intended to illustrate the idealised rate achievable under favourable
geometric conditions and to connect entropy-regularised RL with the Wasserstein
convergence literature. \textbf{No algebraic rate is derived as a fallback under
weaker conditions; this remains an open problem.}
\end{remark}

\begin{table}[ht]
\centering
\small
\caption{Summary of Assumptions and Their Implications}
\label{tab:rl-assumptions}
\begin{tabular}{p{2.5cm}p{2.2cm}p{2.5cm}p{4cm}}
\toprule
\textbf{Assumption} & \textbf{Required for} & \textbf{Practical Relevance} &
    \textbf{Relaxation Possibility} \\
\midrule
Compact $\Theta, \Phi$ & Theorem~\ref{thm:rl-subsequential} & Violated by deep nets &
    Spectral normalisation heuristics \\
Lipschitz approximators & Theorem~\ref{thm:rl-subsequential} & Violated by ReLU nets
    & Smooth activations (tanh, GELU) \\
Stable limiting ODEs & Theorem~\ref{thm:rl-subsequential} & Plausible for LQ, unknown
    for general & Lyapunov analysis \\
Log-Sobolev inequality & Theorem~\ref{thm:rl-exponential} & Only for strongly
    log-concave policies & Poincar\'e inequality may give algebraic rates \\
\bottomrule
\end{tabular}
\end{table}


\section{Proof of Subsequential Almost-Sure Convergence}
\label{sec:rl-proof-subsequential}

The proof adapts \citet{borkar1997,borkar2008}'s (1997, 2008) ODE method for two-timescale stochastic
approximation to the three-timescale, measure-valued setting. The key steps are:

\begin{enumerate}
    \item \textbf{Tightness and compactness}: Show that the sequence of empirical
    measures $\{\hat\mu_n\}$ is tight in $\mathcal{P}_2(\mathbb{R}^d)$ and that the
    parameter sequences $\{\theta_n\}, \{\phi_n\}$ remain in compact sets almost
    surely.
    \item \textbf{Limiting ODE for the critic}: Prove that on the fast timescale,
    with $\phi$ and $\hat\mu$ frozen, the critic parameters $\theta_n$ track the
    solution of the ODE \eqref{eq:critic-ode}, converging to an equilibrium
    $\theta^*(\phi, \hat\mu)$.
    \item \textbf{Limiting ODE for the actor}: On the intermediate timescale, with
    $\hat\mu$ frozen and the critic converged, the actor parameters $\phi_n$ track
    the ODE \eqref{eq:actor-ode}.
    \item \textbf{Limiting SPDE for the measure}: On the slowest timescale, the
    empirical measure $\hat\mu_n$ converges weakly to the solution of the
    Fokker--Planck SPDE \eqref{eq:measure-filter-spde}.
    \item \textbf{Fixed-point argument}: The contraction property of the MFG map
    $\Phi$ (Theorem~\ref{thm:wellposedness}) ensures that the coupled limiting
    dynamics have a unique globally attracting fixed point, which is precisely the
    MFG equilibrium. By the standard stochastic approximation theory, the discrete
    algorithm converges almost surely to this attractor along subsequences.
\end{enumerate}

\subsection{Step 1: Compactness and Tightness}

The boundedness of the function approximators and the Lipschitz properties of the
coefficients ensure that the state process $X_n$ remains in a compact set with high
probability. Specifically, under Assumptions 1.13--1.19, one can show that
$\sup_n \mathbb{E}[|X_n|^2] < \infty$, which implies that the sequence of empirical
measures $\{\hat\mu_n\}$ is tight in $\mathcal{P}_2(\mathbb{R}^d)$. By Prokhorov's
theorem, every subsequence has a further subsequence converging weakly to some limit
measure. The parameters $\theta_n, \phi_n$ are assumed to belong to a compact set
$\Theta \times \Phi$, which can be enforced by projection.

\subsection{Step 2: Critic Convergence (Fast Timescale)}

Fix the actor parameters $\phi$ and the measure flow $\hat\mu$. The critic update is a
standard stochastic approximation for solving the BSDE \eqref{eq:adjoint-bsde}.
Define the vector field $h_\gamma(\theta) = \mathbb{E}_{X \sim \hat\mu}[\delta(\theta,
\phi, \hat\mu) \nabla_\theta p_\theta(X, \hat\mu)]$. By the theory of stochastic
approximation (\citep{borkar2008}, Chapter 2), if $h_\gamma$ is Lipschitz and the ODE
$\dot\theta = h_\gamma(\theta)$ has a globally asymptotically stable equilibrium
$\theta^*$, then $\theta_n \to \theta^*$ almost surely. The equilibrium $\theta^*$
corresponds to the solution of the BSDE, i.e., the true adjoint process under the
given policy and measure.

The global asymptotic stability of the critic ODE follows from the contraction
property of the BSDE (see El Karoui et al., 1997, or \citep{zhang2017}, Chapter 4). Under
Lipschitz conditions on the coefficients, the BSDE solution is unique and depends
continuously on the parameters. The temporal difference error $\delta(\theta, \phi,
\hat\mu)$ is designed so that the ODE drives $\theta$ toward this solution.

\subsection{Step 3: Actor Convergence (Intermediate Timescale)}

With the critic converged to $\theta^*(\phi, \hat\mu)$, the actor update approximates
the ODE $\dot\phi = \mathbb{E}_{X \sim \hat\mu}[\nabla_\phi H(X, \hat\mu, p_{\theta^*},
\pi_\phi) \cdot \nabla_\phi \pi_\phi]$. This is a policy gradient ascent on the
expected Hamiltonian. Under Assumption 3.5 (concavity of $H$ in $\alpha$), the
objective is concave in the policy parameters (for appropriate policy
parametrisations), guaranteeing convergence to the unique maximiser. The timescale
separation $\beta_n/\gamma_n \to 0$ ensures that the critic tracking error vanishes,
so the actor ODE is indeed the correct limiting dynamics.

\subsection{Step 4: Measure Filter Convergence (Slowest Timescale)}

The empirical measure update $\hat\mu_{n+1} = (1-\eta_n)\hat\mu_n + \eta_n
\delta_{X_{n+1}}$ is a stochastic approximation for the conditional law of the state
process. In the limit $\eta_n \to 0$, this converges weakly to the solution of the
Fokker--Planck SPDE \eqref{eq:measure-filter-spde}. The convergence is in the sense of
weak convergence of measure-valued processes; see, for example, the propagation of
chaos results for interacting particle systems (\citep{carmonadelarue2018}, Vol.\ I,
Chapter 1) and the stochastic approximation results for measure-valued processes
(\citep{borkar2008}, Chapter 7). The contraction property of the MFG map $\Phi$
(Theorem~\ref{thm:wellposedness}) guarantees that the limiting measure is unique and
that the discrete approximation converges to it along subsequences.

\subsection{Step 5: Coupling and Subsequential Convergence}

The three limiting dynamics \eqref{eq:critic-ode}, \eqref{eq:actor-ode}, and
\eqref{eq:measure-filter-spde} are coupled. However, due to the timescale separation,
we can analyse them sequentially. The key observation is that the MFG fixed-point map
$\Phi$ is a contraction in the Wasserstein metric. Any fixed point of the coupled
limiting system is precisely an MFG equilibrium. By Theorem~\ref{thm:wellposedness},
this equilibrium is unique. Therefore, the stochastic approximation algorithm, which
asymptotically tracks the limiting ODEs, must visit neighbourhoods of this
equilibrium infinitely often. This yields the $\liminf$ convergence statement.

\paragraph{Why only subsequential convergence?} The ODE method guarantees that the
iterates converge to the internally chain recurrent set of the limiting ODE. In the
presence of multiple attractors or a continuum of equilibria, the iterates may wander
without converging to a single point. Here, the uniqueness of the MFG equilibrium
(Theorem~\ref{thm:wellposedness}) ensures that the only attractor is the equilibrium
itself. However, because the measure-valued component is only shown to converge
weakly, and because the stochastic approximation is subject to persistent noise, we
can only assert that the iterates return to arbitrarily small neighbourhoods of the
equilibrium infinitely often. \textbf{Full almost-sure convergence would require
proving that the sequence is eventually trapped in a neighbourhood, which typically
requires a Lyapunov function with negative drift outside the equilibrium --- a
stronger condition not verified here.}

The subsequential almost-sure convergence statement in
Theorem~\ref{thm:rl-subsequential} follows by combining these steps with the standard
stochastic approximation convergence theorem (\citep{borkar2008}, Theorem 2.1). A detailed
verification of the conditions of \citet{borkar1997,borkar2008}'s theorem for the three-timescale,
measure-valued setting is provided in Appendix~A.


\section{Exponential Convergence under Entropy Regularisation}
\label{sec:rl-entropy}

We now prove that adding entropy regularisation leads to an exponential convergence
rate under the log-Sobolev assumption.

\subsection{Entropy-Regularised Objective}

The entropy-regularised objective for the representative agent is
\[
    J_\beta(\alpha; \mu) = \mathbb{E}\left[ \int_0^T \big( f(t,X_t,\mu_t,\alpha_t) -
    \beta \log \pi_t(\alpha_t \mid X_t, \mu_t) \big)\, dt + g(X_T,\mu_T) \right],
\]
where $\pi_t(\cdot \mid X_t, \mu_t)$ is the policy density with respect to a base
measure on $A$. The corresponding entropy-regularised Hamiltonian is
\[
    H_\beta(t,x,\mu,p,\alpha) = p \cdot b(t,x,\mu,\alpha) + f(t,x,\mu,\alpha) - \beta
    \log \pi(\alpha \mid x, \mu).
\]
The optimal policy satisfies a Gibbs distribution:
\begin{equation}
    \pi^*(\alpha \mid x, \mu) \propto \exp\left( \frac{1}{\beta} \big( p \cdot
    b(t,x,\mu,\alpha) + f(t,x,\mu,\alpha) \big) \right).
    \label{eq:gibbs-policy}
\end{equation}

\subsection{Log-Sobolev Inequality and Exponential Convergence}

The log-Sobolev inequality (LSI) for the policy class states that for any two
policies $\pi, \pi'$,
\[
    \mathrm{KL}(\pi \| \pi') \leq \frac{C_{\mathrm{Lip}}}{2} \mathbb{E}_{X \sim \mu}
    \big[ \|\nabla_\phi \log \pi_\phi(X) - \nabla_{\phi'} \log \pi_{\phi'}(X)\|^2
    \big],
\]
where $\phi, \phi'$ are the parameters of $\pi, \pi'$ respectively. This inequality
implies that the policy gradient dynamics converge exponentially fast to the optimal
policy. The connection between LSI and exponential convergence of Langevin dynamics
and policy gradient methods is well established (see, e.g., \citep{malizhang2024}.

Under the LSI, the continuous-time actor dynamics satisfy
\[
    \frac{d}{dt} \mathrm{KL}(\pi_{\phi_t} \| \pi^*) \leq -2\lambda\,
    \mathrm{KL}(\pi_{\phi_t} \| \pi^*),
\]
with $\lambda = \beta/(2C_{\mathrm{Lip}})$. By the Talagrand transportation
inequality, which relates KL divergence to the Wasserstein distance (specifically,
$W_2(\mu,\nu)^2 \leq 2 C_{\mathrm{Lip}}\, \mathrm{KL}(\mu\|\nu)$ for measures
satisfying LSI), this translates to
\[
    \frac{d}{dt} W_2(\mu_t, \mu^*)^2 \leq -2\lambda W_2(\mu_t, \mu^*)^2.
\]
Applying Gronwall's inequality yields the exponential convergence claimed in
Theorem~\ref{thm:rl-exponential}.

\subsection{Discussion of the Log-Sobolev Assumption}

The log-Sobolev inequality is a strong structural condition. It is satisfied, for
example, by Gaussian policies or policies that are strongly log-concave in the
parameters. In the context of neural network policies, the LSI typically fails
unless explicit regularisation (such as weight decay or spectral normalisation) is
imposed. Recent work on Lata\l a--Oleszkiewicz inequalities and non-convex
isoperimetry (see Chapter 10) offers potential pathways to relax this assumption. The
exponential rate in Theorem~\ref{thm:rl-exponential} provides a theoretical upper
bound on what can be achieved under idealised conditions.


\section{Practical Implementation Considerations and Limitations}
\label{sec:rl-implementation}

\subsection{Function Approximation}

In practice, the critic $p_\theta$ and the actor $\pi_\phi$ are implemented as
neural networks. To approximate the theoretical assumptions, one can use bounded
activation functions (e.g., tanh) and constrain the weight matrices to have bounded
spectral norm (e.g., via spectral normalisation). The parameter spaces are then
compact, and the networks are Lipschitz continuous. While this does not fully
satisfy the compactness assumption (which requires a finite-dimensional compact set),
it provides a practical heuristic.

\subsection{Computing the Wasserstein Distance}

The critic and actor updates require evaluating the measure flow $\hat\mu_t$, which
is an empirical distribution of particles. The Wasserstein distance
$W_2(\hat\mu_t, \hat\mu_s)$ or gradients with respect to the measure can be
approximated using the Sinkhorn algorithm (\Cref{sec:mp-sinkhorn}) or, in one
dimension, via the quantile formula \eqref{eq:1d-wasserstein}. For high-dimensional
problems, sliced Wasserstein distances or maximum mean discrepancies (MMD) can be
used as computationally cheaper alternatives, though they do not enjoy the same
theoretical guarantees.

\subsection{Parallelisation and Scalability}

The algorithm is naturally parallelisable: multiple agents can be simulated in
parallel, each maintaining its own critic and actor, while sharing the population
measure $\hat\mu_t$ via a centralised parameter server. The common noise $W^0$ must
be identical across all agents to preserve the synchronous coupling required for the
contraction property.

\subsection{Learning Rate Schedules}

In practice, the Robbins--Monro conditions are approximated by using constant
learning rates with occasional decay. The timescale separation is enforced by
choosing $\gamma \gg \beta \gg \eta$. Typical values are
$\gamma = 10^{-3}, \beta = 10^{-4}, \eta = 10^{-5}$. \textbf{However, these fixed
ratios do not satisfy the asymptotic condition $\beta_t/\gamma_t \to 0$; the
finite-time behaviour may exhibit transient oscillations not captured by the
theory.}

\subsection{Limitations of the Theoretical Guarantees}

The convergence guarantees in Theorems \ref{thm:rl-subsequential} and
\ref{thm:rl-exponential} are subject to the following important caveats:

\begin{itemize}
    \item \textbf{Compactness and Lipschitz assumptions}: Not satisfied by standard
    deep neural networks.
    \item \textbf{Subsequential convergence}: Only $\liminf W_2 = 0$ is guaranteed,
    not full convergence.
    \item \textbf{Idealised LSI}: Exponential rate only under strong log-concavity,
    unlikely to hold in practice.
    \item \textbf{Asymptotic timescale separation}: Fixed learning rate ratios
    violate the theoretical condition.
\end{itemize}

These limitations do not invalidate the theoretical contributions, but they
underscore the gap between theory and practice. The experimental results in Chapter 6
demonstrate that, despite these theoretical gaps, the proposed architecture achieves
strong empirical performance.


\section{Conclusion and Connection to Chapter 6}
\label{sec:rl-conclusion}

We have established convergence guarantees for a three-timescale actor-critic
algorithm in the common-noise MFG setting. The subsequential almost-sure convergence
(Theorem~\ref{thm:rl-subsequential}) relies on the contraction property of the MFG
map (Theorem~\ref{thm:wellposedness}) and the ODE method for stochastic
approximation. The exponential rate (Theorem~\ref{thm:rl-exponential}) under entropy
regularisation and log-Sobolev inequality provides an idealised benchmark for fast
convergence.

These theoretical results form the foundation for the hierarchical architecture
developed in Chapter 6. The three natural timescales --- idiosyncratic diffusion,
mean-field interaction, and common-noise drift --- map directly to the critic,
actor, and meta-learner components of a hierarchical RL system. In Chapter 6, we
formalise this correspondence and demonstrate how state-of-the-art algorithms (SAC,
PPO, TD3) can be orchestrated within this framework to achieve efficient learning in
practice.
